\subsection{The arbitrary-graphon comparison}
\label{sec:graphon-comparison}

We now restore the remainder $E$ and convert the lower bound for the auxiliary
Lagrangian difference into the corresponding bound for the Lagrangian
difference between the actual graphons $W$ and $W_h$.
By \Cref{lem:localization-rank-one}, every feasible graphon satisfying
$I_{p_h}(W)\le I_{p_h}(W_h)$ can be written as $W=f\otimes f+E$, while
\Cref{lem:auxiliary-lagrangian-bound} controls the auxiliary Lagrangian
difference between $f$ and $f_h$.  The next lemma estimates the contribution
of $E$ and gives the required full-graphon bound.

\begin{lemma}[Lower bound for the full-graphon Lagrangian difference]
\label{lem:graphon-lagrangian-bound}
There exist $\rho_0>0$ and $C_{d,m}\ge1$ such that for
every $0<\rho<\rho_0$, one can choose $C_{d,\rho}\ge1$ and $h_\rho>0$ with
the following property.  Let $0<h<h_\rho$, and let $W$ satisfy the hypotheses of
\Cref{lem:localization-rank-one}.  Write $W=f\otimes f+E$ for the
decomposition provided there, and let $\nu$ be the distribution of the
values of $f$.  Then,
\begin{equation}\label{eq:graphon-lagrangian-bound}
\begin{aligned}
  &\left[I_{p_h}(W)-\mu_h\{t(H,W)-r_h^m\}\right]
   -\left[I_{p_h}(W_h)-\mu_h\{t(H,W_h)-r_h^m\}\right]\\
  &\qquad\qquad\qquad\qquad\qquad\qquad\qquad=I_{p_h}(W)-I_{p_h}(W_h)-\mu_h\{t(H,W)-r_h^m\}\\
  &\qquad\qquad\qquad\qquad\qquad\qquad\qquad
  \ge C_{d,\rho}^{-1}\!\left[A_{h,\rho}(\nu)+\nu(\mathcal T_\rho)+\|E\|_2^2\right]
  -C_{d,m}\Delta_h(\nu)^2.
\end{aligned}
\end{equation}
\end{lemma}

\begin{proof}
Put $q:=q(f)=m_d(\nu)$ and
$\Omega_\rho:=\{x:f(x)\in\mathcal N_\rho\}$.  Then
$\lambda(\Omega_\rho^c)=\nu(\mathcal T_\rho)$.
Combining the localization results in \Cref{lem:localization-rank-one}
and \eqref{eq:factor-tail-mass}, we may use throughout the proof
\begin{equation}\label{eq:graphon-comparison-estimates}
  \begin{gathered}
    \nu(\mathcal T_\rho)\le C_{d,\rho}h^4,
    \qquad |\Delta_h(\nu)|\le C_dh^2,
    \qquad \|E\|_2\le C_dh,
    \qquad \|f-u_*\|_4\le C_dh,\\
    |t(H,W)-q^v|\le C_{d,m}\|E\|_2^3,
    \qquad \lambda(\Omega_\rho)\ge\frac12,
    \qquad 0\le f\le2,
    \qquad \|E\|_\infty\le5.
  \end{gathered}
\end{equation}
The bound on $\lambda(\Omega_\rho)$ follows from
$\lambda(\Omega_\rho^c)=\nu(\mathcal T_\rho)$, while the last bound uses
$0\le W\le1$ and $0\le f\le2$.
Since
\[
  \mathcal J_h(\nu)
  =\iint\widetilde J_{p_h}(f(x)f(y))\,dx\,dy,
  \qquad
  \mathcal J_h(\nu_h)=I_{p_h}(W_h),
\]
and $t(H,W_h)=q_h^v=r_h^m$, the full-graphon Lagrangian difference has the
exact decomposition
\begin{equation*}
\begin{aligned}
  &I_{p_h}(W)-I_{p_h}(W_h)-\mu_h\{t(H,W)-r_h^m\}\\
  &\quad=\bigl[\mathcal L_h(\nu)-\mathcal L_h(\nu_h)\bigr]
  +\iint\{J_{p_h}(W(x,y))-\widetilde J_{p_h}(f(x)f(y))\}\,dx\,dy\\
  &\qquad-\mu_h\{t(H,W)-q^v\}.
\end{aligned}
\end{equation*}
The first term is controlled by \Cref{lem:auxiliary-lagrangian-bound}, and
the last by the homomorphism-density estimate in
\eqref{eq:rank-one-reduction-bounds}.  It remains to bound the middle
double integral over its central, mixed, and tail--tail parts.

We first consider the central square $\Omega_\rho^2$.  Decrease $\rho_0$ so
that $f(x)f(y)$ lies in a fixed closed subinterval of $(0,1)$ whenever
$x,y\in\Omega_\rho$.  Put
\[
  \|E\|_{2,G}
  :=\left( \iint_G E(x,y)^2\,dx\,dy \right)^{1/2}
  \qquad (G \subset[0,1]^2).
\]
On $\Omega_\rho^2$, the identity $\mathcal M_h'(xy)=J_{p_h}'(xy)-\gamma_h(xy)^{d-1}$ yields
the decomposition of the integrand:
\begin{equation*}
  J_{p_h}(W)-J_{p_h}(f\otimes f)
  = \bigl[ J_{p_h}(W)-J_{p_h}(f\otimes f) -(J_{p_h})'(f\otimes f)E \bigr]
  +\mathcal M_h'(f\otimes f)E
      +\gamma_h(f\otimes f)^{d-1}E.
\end{equation*}
In this step, we bound the integrals of these three terms on $\Omega_\rho^2$ separately.
Since $J_{p_h}''(z)=1/(z(1-z))\ge4$, strong convexity gives the first bound
\[
  \iint_{\Omega_\rho^2}
  \{J_{p_h}(W)-J_{p_h}(f\otimes f)-(J_{p_h})'(f\otimes f)E\}
  \ge2\|E\|_{2,\Omega_\rho^2}^2.
\]
For the second bound, we divide $\mathcal M_h'$ by $Q_h$ in each variable and write
\[
  \mathcal M_h'(xy)=Q_h(x)U_h(x,y)+Q_h(y)V_h(x,y)+C_h(x,y),
\]
where $C_h$ has degree at most one in each variable.  The factorization
\eqref{eq:mh-derivative-factorization} then gives
$\mathcal M_h'(xy)=0$ for every $(x,y)\in\{s_h,t_h\}^2$.  Since
$Q_h(s_h)=Q_h(t_h)=0$, the bilinear remainder $C_h$ must vanish identically.
At $h=0$, $\mathcal M_0'(xy)$ is an analytic multiple of $(xy-r_*)^3$ and
$Q_0(z)=(z-u_*)^2$, so $U_0$ and $V_0$ vanish at $(u_*,u_*)$.  By
joint continuity, after decreasing $\rho_0$ and $h_\rho$, we may ensure that
$|U_h|$ and $|V_h|$ are at most $\sqrt{a/4}$ on $\mathcal N_\rho^2$, and
\[
  |\mathcal M_h'(xy)|
  \le\sqrt{\frac a4}
    \{|Q_h(x)|+|Q_h(y)|\}
  \qquad(x,y\in\mathcal N_\rho,\ 0<h<h_\rho).
\]
Then, $uv \le (1/4)u^2+v^2$ and
$(u+v)^2\le 2(u^2+v^2)$ give
\begin{align*}
  \left|\iint_{\Omega_\rho^2} \{\mathcal M_h'(f(x)f(y))\} E(x,y)\,dx\,dy\right|
  &\le \frac{1}{4} \iint_{\Omega_\rho^2} \frac{a}{4}\{|Q_h(f(x))|+|Q_h(f(y))|\}^2 \,dx\,dy + \|E\|_{2,\Omega_\rho^2}^2 \\
  &\le \frac{a}{16} \iint_{\Omega_\rho^2} 2|Q_h(f(x))|^2+2|Q_h(f(y))|^2 \,dx\,dy + \|E\|_{2,\Omega_\rho^2}^2 \\
  &\le \frac a4 A_{h,\rho}(\nu) + \|E\|_{2,\Omega_\rho^2}^2.
\end{align*}
For the third bound, the orthogonality in \eqref{eq:rank-one-orthogonality} moves
the remaining linear term to the tail--tail square:
\[
  \left| \iint_{\Omega_\rho^2}\gamma_h(f(x)f(y))^{d-1}E(x,y)\,dx\,dy \right|
  = \left| \iint_{(\Omega_\rho^c)^2}\gamma_h(f(x)f(y))^{d-1}E(x,y)\,dx\,dy \right|
  \le C_d \nu(\mathcal T_\rho)^2.
\]
Combining the three bounds gives
\begin{equation*}
  \iint_{\Omega_\rho^2} \{J_{p_h}(W)-J_{p_h}(f\otimes f)\}
  \ge \|E\|_{2,\Omega_\rho^2}^2 -\frac a4 A_{h,\rho}(\nu)-C_d\nu(\mathcal T_\rho)^2.
\end{equation*}

We next consider the mixed rectangles.  For almost every $x$, the
orthogonality in
\eqref{eq:rank-one-orthogonality}, together with
\eqref{eq:graphon-comparison-estimates}, gives
\begin{equation*}
  u_*^{d-1}\left|\int_{\Omega_\rho}E(x,y)\,dy\right|
  \le\int_{\Omega_\rho}
    |f(y)^{d-1}-u_*^{d-1}|\,|E(x,y)|\,dy
  +\int_{\Omega_\rho^c}f(y)^{d-1}|E(x,y)|\,dy
  \le C_{d,\rho}h.
\end{equation*}
Consequently, if $\overline W_x:= (\int_{\Omega_\rho}W(x,y)\,dy)/\lambda(\Omega_\rho)$,
then
\begin{equation*}
  |\overline W_x-u_*f(x)|
  \le |f(x)|\left|
    \frac1{\lambda(\Omega_\rho)}\int_{\Omega_\rho} (f(y)-u_*)\,dy\right|
    +\frac1{\lambda(\Omega_\rho)}\left|\int_{\Omega_\rho}E(x,y)\,dy\right|
  \le C_{d,\rho}h.
\end{equation*}
Since $\overline W_x\in[0,1]$, we have
$J_{p_h}(\overline W_x)=\widetilde J_{p_h}(\overline W_x)$.
By Jensen's and H\"older's inequalities, with
\eqref{eq:continuation-kernel-bounds} and
\eqref{eq:graphon-comparison-estimates},
\begin{equation*}
\begin{aligned}
&\int_{\Omega_\rho}
  \{J_{p_h}(W(x,y))-\widetilde J_{p_h}(f(x)f(y))\}\,dy\\
&\quad\ge \lambda(\Omega_\rho) \left\{J_{p_h}(\overline W_x)-\widetilde J_{p_h}(u_*f(x))\right\}
-\int_{\Omega_\rho} \left\{\widetilde J_{p_h}(f(x)f(y))-\widetilde J_{p_h}(u_*f(x))\right\}\,dy\\
&\quad\ge-C_d |\overline W_x - u_*f(x)|^{1/2} -C_d\int_{\Omega_\rho}|f(y)-u_*|^{1/2}\,dy \ge-C_{d,\rho}h^{1/2} \ge -\frac{b_\rho}{16}.
\end{aligned}
\end{equation*}
Integrating this bound over $x\in\Omega_\rho^c$ yields
\begin{equation*}
\begin{aligned}
  2\iint_{\Omega_\rho^c\times \Omega_\rho}
  \{J_{p_h}(W)-\widetilde J_{p_h}(f\otimes f)\}
  \ge -\frac{b_\rho}{8}\nu(\mathcal T_\rho).
\end{aligned}
\end{equation*}

It remains to consider the tail--tail square.  The uniform bound in
\eqref{eq:continuation-kernel-bounds} gives
\[
  \iint_{(\Omega_\rho^c)^2}
  \{J_{p_h}(W)-\widetilde J_{p_h}(f\otimes f)\}
  \ge-C_d\nu(\mathcal T_\rho)^2.
\]

Combining all bounds with \eqref{eq:auxiliary-lagrangian-bound}
and $|t(H,W)-q^v|\le C_{d,m}\|E\|_2^3$ gives
\begin{equation*}
\begin{aligned}
  I_{p_h}(W)-I_{p_h}(W_h)-\mu_h\{t(H,W)-r_h^m\}
  &\ge\frac a4 A_{h,\rho}(\nu)
    +\frac{3b_\rho}{8}\nu(\mathcal T_\rho)
    +\|E\|_{2,\Omega_\rho^2}^2\\
  &\quad-C_{d,m}\Delta_h(\nu)^2
    -C_d\nu(\mathcal T_\rho)^2
    -C_{d,m}\|E\|_2^3.
\end{aligned}
\end{equation*}
Since $\|E\|_\infty\le5$ and
$\lambda([0,1]^2\setminus\Omega_\rho^2)\le2\nu(\mathcal T_\rho)$, choose
$C_{d,\rho}\ge1$ so that $50C_{d,\rho}^{-1}\le b_\rho/16$.  Then
\begin{equation*}
  \|E\|_{2,\Omega_\rho^2}^2
  \ge C_{d,\rho}^{-1}\|E\|_{2,\Omega_\rho^2}^2
  =C_{d,\rho}^{-1}\left(
    \|E\|_2^2-\|E\|_{2,[0,1]^2\setminus\Omega_\rho^2}^2\right)
  \ge C_{d,\rho}^{-1}\|E\|_2^2
    -50C_{d,\rho}^{-1}\nu(\mathcal T_\rho)
  \ge C_{d,\rho}^{-1}\|E\|_2^2
    -\frac{b_\rho}{16}\nu(\mathcal T_\rho).
\end{equation*}
By \eqref{eq:graphon-comparison-estimates}, we may also decrease $h_\rho$ so that
$C_d\nu(\mathcal T_\rho)\le b_\rho/16$ and $C_{d,m}\|E\|_2\le C_{d,\rho}^{-1}/2$.
Consequently,
\begin{equation*}
  I_{p_h}(W)-I_{p_h}(W_h)-\mu_h\{t(H,W)-r_h^m\}
  \ge\frac a4 A_{h,\rho}(\nu)
  +\frac{b_\rho}{4}\nu(\mathcal T_\rho)
  +\frac{1}{2}C_{d,\rho}^{-1}\|E\|_2^2
  -C_{d,m}\Delta_h(\nu)^2.
\end{equation*}
Increasing $C_{d,\rho}$ if necessary proves \eqref{eq:graphon-lagrangian-bound}.
\end{proof}
